\documentclass[10pt,a4paper]{amsart}
\usepackage{amsmath,amssymb,amsthm,graphicx}
\usepackage{verbatim}
\newenvironment{reference}{\comment}{\endcomment}
%\newenvironment{reference}{}{}
\newenvironment{slogan}{\comment}{\endcomment}
\newenvironment{history}{\comment}{\endcomment}
% For commutative diagrams we use Xy-pic
\usepackage[all]{xy}

% We use 2cell for 2-commutative diagrams.
\xyoption{2cell}
% Theorem environments.
%
\theoremstyle{plain}
\newtheorem{theorem}[subsection]{Theorem}
\newtheorem{proposition}[subsection]{Proposition}
\newtheorem{lemma}[subsection]{Lemma}

\theoremstyle{definition}
\newtheorem{definition}[subsection]{Definition}
\newtheorem{example}[subsection]{Example}
\newtheorem{exercise}[subsection]{Exercise}
\newtheorem{situation}[subsection]{Situation}

\theoremstyle{remark}
\newtheorem{remark}[subsection]{Remark}
\newtheorem{remarks}[subsection]{Remarks}

\numberwithin{equation}{subsection}

% Macros
%
\def\lim{\mathop{\mathrm{lim}}\nolimits}
\def\colim{\mathop{\mathrm{colim}}\nolimits}
\def\Spec{\mathop{\mathrm{Spec}}}
\def\Hom{\mathop{\mathrm{Hom}}\nolimits}
\def\Ext{\mathop{\mathrm{Ext}}\nolimits}
\def\SheafHom{\mathop{\mathcal{H}\!\mathit{om}}\nolimits}
\def\SheafExt{\mathop{\mathcal{E}\!\mathit{xt}}\nolimits}
\def\Sch{\mathit{Sch}}
\def\Mor{\mathop{\mathrm{Mor}}\nolimits}
\def\Ob{\mathop{\mathrm{Ob}}\nolimits}
\def\Sh{\mathop{\mathit{Sh}}\nolimits}
\def\NL{\mathop{N\!L}\nolimits}
\def\CH{\mathop{\mathrm{CH}}\nolimits}
\def\proetale{{pro\text{-}\acute{e}tale}}
\def\etale{{\acute{e}tale}}
\def\QCoh{\mathit{QCoh}}
\def\Ker{\mathop{\mathrm{Ker}}}
\def\Im{\mathop{\mathrm{Im}}}
\def\Coker{\mathop{\mathrm{Coker}}}
\def\Coim{\mathop{\mathrm{Coim}}}

% Boxtimes
%
\DeclareMathSymbol{\boxtimes}{\mathbin}{AMSa}{"02}

%
% Macros for moduli stacks/spaces
%
\def\QCohstack{\mathcal{QC}\!\mathit{oh}}
\def\Cohstack{\mathcal{C}\!\mathit{oh}}
\def\Spacesstack{\mathcal{S}\!\mathit{paces}}
\def\Quotfunctor{\mathrm{Quot}}
\def\Hilbfunctor{\mathrm{Hilb}}
\def\Curvesstack{\mathcal{C}\!\mathit{urves}}
\def\Polarizedstack{\mathcal{P}\!\mathit{olarized}}
\def\Complexesstack{\mathcal{C}\!\mathit{omplexes}}
% \Pic is the operator that assigns to X its picard group, usage \Pic(X)
% \Picardstack_{X/B} denotes the Picard stack of X over B
% \Picardfunctor_{X/B} denotes the Picard functor of X over B
\def\Pic{\mathop{\mathrm{Pic}}\nolimits}
\def\Picardstack{\mathcal{P}\!\mathit{ic}}
\def\Picardfunctor{\mathrm{Pic}}
\def\Deformationcategory{\mathcal{D}\!\mathit{ef}}
\usepackage{cleveref}
\crefname{originallemma}{original lemma}{original lemmas}
\Crefname{originallemma}{Original lemma}{Original lemmas}
\makeatletter
\newlabel{lemma-base-change-flat}{{lemma-base-change-flat}{source}{}{sourcepointer-1}{}}
\newlabel{lemma-base-change-flat@cref}{{[originallemma][0][]lemma-base-change-flat}{[1][0][]source}{}{}{}}
\makeatother
\crefname{originallemma}{original lemma}{original lemmas}
\Crefname{originallemma}{Original lemma}{Original lemmas}
\makeatletter
\newlabel{lemma-composition-flat}{{lemma-composition-flat}{source}{}{sourcepointer-2}{}}
\newlabel{lemma-composition-flat@cref}{{[originallemma][0][]lemma-composition-flat}{[1][0][]source}{}{}{}}
\makeatother
\crefname{originalreference}{original reference}{original references}
\Crefname{originalreference}{Original reference}{Original references}
\makeatletter
\newlabel{spaces-morphisms-lemma-base-change-module-flat}{{spaces-morphisms-lemma-base-change-module-flat}{source}{}{sourcepointer-3}{}}
\newlabel{spaces-morphisms-lemma-base-change-module-flat@cref}{{[originalreference][0][]spaces-morphisms-lemma-base-change-module-flat}{[1][0][]source}{}{}{}}
\makeatother
\crefname{originalreference}{original reference}{original references}
\Crefname{originalreference}{Original reference}{Original references}
\makeatletter
\newlabel{spaces-morphisms-lemma-smooth-flat}{{spaces-morphisms-lemma-smooth-flat}{source}{}{sourcepointer-4}{}}
\newlabel{spaces-morphisms-lemma-smooth-flat@cref}{{[originalreference][0][]spaces-morphisms-lemma-smooth-flat}{[1][0][]source}{}{}{}}
\makeatother
\crefname{originalreference}{original reference}{original references}
\Crefname{originalreference}{Original reference}{Original references}
\makeatletter
\newlabel{stacks-properties-lemma-composition-surjective}{{stacks-properties-lemma-composition-surjective}{source}{}{sourcepointer-5}{}}
\newlabel{stacks-properties-lemma-composition-surjective@cref}{{[originalreference][0][]stacks-properties-lemma-composition-surjective}{[1][0][]source}{}{}{}}
\makeatother
\title{Flatness of a morphism of algebraic stacks is checked after a flat surjective base change}
\author{The Stacks Project Authors}
\date{stacks-project:stacks-morphisms.tex:local-snapshot-20261012; GNU FDL 1.2 or later}
\begin{document}
\maketitle
\noindent\textit{Editorial note.} All delivered mathematical text is the author's own, copied from the pinned Stacks Project chapter source sources/stacks-algebraic-stacks-morphisms/pinned-original-source.tex; no statement or proof step was rewritten, reordered or removed.\par\medskip
\noindent\textit{Editorial note.} Preamble provenance: the Stacks Project keeps one single shared preamble for every chapter; this card takes it directly from the staged sources/stacks-algebraic-stacks-morphisms/preamble.tex as three verbatim spans (lines 8-12, 14-18 and 154-217), so no derived file was written and nothing was added or redefined. The wrapper class and the base packages are supplied by the assembler.\par\medskip
\noindent\textit{Editorial note.} Bibliography provenance: neither the counted statement nor its proof contains a citation (this lane rescans the whole span before writing the card), so the card carries no bibliography and no entry was reconstructed.\par\medskip
\noindent\textit{Editorial note.} Selection provenance: this card was opened under the file-based control contract (CONTROL.md), with the mechanical pre-card gate run on its statement and proof interval before the card was written.\par\medskip
\noindent\textit{Allowed standard prerequisites (editorial).}\begin{itemize}
\item The statement that a smooth morphism of algebraic spaces is flat, and the composition permanence of flatness. Invoked at the start of the counted proof (source lines 5398-5403) as cross-chapter pointers. Makes the composite presentation morphism flat and surjective.
\item The base change formula for flatness. Invoked inside the counted proof (source lines 5404-5406). Justifies replacing the test object by a smooth presentation of it.
\item The definition of flatness for a morphism of algebraic stacks through presentations. Used inside the counted proof (source lines 5415-5420). Turns the assertion into a statement about the presenting morphisms of algebraic spaces.
\item The statement that flatness of a morphism of algebraic spaces descends along a flat surjective base change. Invoked at the end of the counted proof (source lines 5425-5428) as a cross-chapter pointer. Concludes flatness of the presenting morphism before the base change.
\end{itemize}\medskip
\ifdefined\hypertarget\hypertarget{sourcepointer-1}{}\fi
\noindent\textit{Source reference.} Original lemma lemma-base-change-flat: sources/stacks-algebraic-stacks-morphisms/pinned-original-source.tex, source lines 5376--5376. Navigation pointer only: the counted text refers to the Stacks item lemma-base-change-flat, which is an existing labeled result of this same chapter and is not part of this card. The pointer carries the target's own label as its number and its definition line as its span; its content is not reproduced here.\par\medskip
\ifdefined\hypertarget\hypertarget{sourcepointer-2}{}\fi
\noindent\textit{Source reference.} Original lemma lemma-composition-flat: sources/stacks-algebraic-stacks-morphisms/pinned-original-source.tex, source lines 5363--5363. Navigation pointer only: the counted text refers to the Stacks item lemma-composition-flat, which is an existing labeled result of this same chapter and is not part of this card. The pointer carries the target's own label as its number and its definition line as its span; its content is not reproduced here.\par\medskip
\ifdefined\hypertarget\hypertarget{sourcepointer-3}{}\fi
\noindent\textit{Source reference.} Original reference spaces-morphisms-lemma-base-change-module-flat: sources/stacks-algebraic-stacks-morphisms/pinned-original-source.tex, source lines 5428--5428. Navigation pointer only: the counted text refers to the Stacks item spaces-morphisms-lemma-base-change-module-flat, which belongs to another chapter of the project and is not part of this card. The pointer carries the target's own label as its number because this chapter's staged snapshot does not contain the other chapter's tag table; its content is not reproduced here.\par\medskip
\ifdefined\hypertarget\hypertarget{sourcepointer-4}{}\fi
\noindent\textit{Source reference.} Original reference spaces-morphisms-lemma-smooth-flat: sources/stacks-algebraic-stacks-morphisms/pinned-original-source.tex, source lines 5400--5400. Navigation pointer only: the counted text refers to the Stacks item spaces-morphisms-lemma-smooth-flat, which belongs to another chapter of the project and is not part of this card. The pointer carries the target's own label as its number because this chapter's staged snapshot does not contain the other chapter's tag table; its content is not reproduced here.\par\medskip
\ifdefined\hypertarget\hypertarget{sourcepointer-5}{}\fi
\noindent\textit{Source reference.} Original reference stacks-properties-lemma-composition-surjective: sources/stacks-algebraic-stacks-morphisms/pinned-original-source.tex, source lines 5403--5403. Navigation pointer only: the counted text refers to the Stacks item stacks-properties-lemma-composition-surjective, which belongs to another chapter of the project and is not part of this card. The pointer carries the target's own label as its number because this chapter's staged snapshot does not contain the other chapter's tag table; its content is not reproduced here.\par\medskip
\begin{lemma}
\label{lemma-descent-flat}
Let $f : \mathcal{X} \to \mathcal{Y}$ be a morphism of algebraic stacks.
Let $\mathcal{Z} \to \mathcal{Y}$ be a surjective flat morphism of algebraic
stacks. If the base change
$\mathcal{Z} \times_\mathcal{Y} \mathcal{X} \to \mathcal{Z}$
is flat, then $f$ is flat.
\end{lemma}

\begin{proof}
Choose an algebraic space $W$ and a surjective smooth morphism
$W \to \mathcal{Z}$. Then $W \to \mathcal{Z}$ is surjective and flat
(Morphisms of Spaces, Lemma \ref{spaces-morphisms-lemma-smooth-flat})
hence $W \to \mathcal{Y}$ is surjective and flat (by
Properties of Stacks, Lemma
\ref{stacks-properties-lemma-composition-surjective}
and
Lemma \ref{lemma-composition-flat}).
Since the base change of
$\mathcal{Z} \times_\mathcal{Y} \mathcal{X} \to \mathcal{Z}$
by $W \to \mathcal{Z}$ is a flat morphism
(Lemma \ref{lemma-base-change-flat})
we may replace $\mathcal{Z}$ by $W$.

\medskip\noindent
Choose an algebraic space $V$ and a surjective smooth morphism
$V \to \mathcal{Y}$. Choose an algebraic space $U$ and a surjective
smooth morphism $U \to V \times_\mathcal{Y} \mathcal{X}$.
We have to show that $U \to V$ is flat. Now we base change
everything by $W \to \mathcal{Y}$: Set $U' = W \times_\mathcal{Y} U$,
$V' = W \times_\mathcal{Y} V$,
$\mathcal{X}' = W \times_\mathcal{Y} \mathcal{X}$,
and $\mathcal{Y}' = W \times_\mathcal{Y} \mathcal{Y} = W$.
Then it is still true that $U' \to V' \times_{\mathcal{Y}'} \mathcal{X}'$
is smooth by base change. Hence by our definition of flat morphisms
of algebraic stacks and the assumption that $\mathcal{X}' \to \mathcal{Y}'$
is flat, we see that $U' \to V'$ is flat. Then, since
$V' \to V$ is surjective as a base change of $W \to \mathcal{Y}$ we see
that $U \to V$ is flat by
Morphisms of Spaces, Lemma
\ref{spaces-morphisms-lemma-base-change-module-flat} (2)
and we win.
\end{proof}

\end{document}
